%%%% PROCESAR con PdfLaTeX !!!!!


\documentclass[12pt]{book}
\usepackage{geometry}\geometry{top=2cm,bottom=2cm,left=3cm,right=3cm}
\usepackage{amssymb}
\usepackage{amsmath}
\usepackage{graphicx}
\usepackage{txfonts}
%\usepackage{hyperref}
\usepackage[hidelinks]{hyperref}
\usepackage[spanish]{babel}
\setcounter{tocdepth}{3}
\usepackage[usenames]{color}

\usepackage[hidelinks]{hyperref}
\hypersetup{
    colorlinks=true,
    linkcolor=black,
    filecolor=magenta,      
    urlcolor=blue,
    pdftitle={Sharelatex Example},
    bookmarks=true,
%    pdfpagemode=FullScreen,
}


\begin{document}
\thispagestyle{empty}

\begin {center}


\includegraphics[scale=.4]{../img/descarga.jpeg}

\medskip
\textbf{\large UNIVERSIDAD DE BUENOS AIRES}

Facultad de Derecho

Departamento: Derecho Privado
\\
Catedra: Oscar Ameal
\\
COMISIONES 6820 y 6821
\vspace{3cm}


\textbf{\large Elementos de Derecho Civil (Parte General)}

\vspace{2cm}

Este es un modesto aporte para los alumnos de la f\'acultad de derecho de la UBA para la carrera de Abogac\'ia.
De ninguna man\'era pretende ser una gu\'ia de estudio, ni remplaza las clases presenciales, el material oficial de la catedra esta disponible en el web site de la m\'ateria.
\\
\url{https://sites.google.com/site/comision606/}

\end {center}


\vspace{2.5cm}

\noindent Autor:\,	Isaac Edgar Camacho Ocampo
 
\noindent Carrera:\,	Abogac\'ia

\vspace{1cm}

\vspace{1cm}

\noindent Buenos Aires, 2020
\vspace{1cm}
\\
\textit{Si se encuentra alg\'un error u omisi\'on en este res\'umen por favor colaborar en \\
\url{https://github.com/IsaacEdgarCamacho/Apuntes/tree/master/Derecho} \quad \\ o escribirme a \\ 
\url{camachoocampo435@est.derecho.uba.ar}
}
\newpage


\tableofcontents
\chapter*{RESUMEN} % si no queremos que añada la palabra "Capitulo"
\addcontentsline{toc}{section}{Resumen} % si queremos que aparezca en el índice
\markboth{RESUMEN}{RESUMEN} % encabezado

Los contenidos que conforman la asignatura corresponden a una Parte General del Derecho Civil, siendo por tal razón, base de sustentación técnico-científica del resto de los conocimientos que el alumno ha de ir incorporando en el devenir de sus estudios universitarios.
\\ \\
El alumno se encuentra en el umbral de su carrera y por lo tanto, su formación y capacitación debe pensarse como una permanente e ininterrumpida aprehensión e interna lización, tanto de conocimientos como de valores y vivencias \\ \\\
En este sentido, los contenidos que conforman el programa de esta materia, se han diseñado en vista de apuntalar al estudiante, en las nociones básicas del Derecho Civil, a partir del estudio y la reflexión sobre sus instituciones fundantes. \\ \\
En este contexto, las expectativas se centran en las posibilidades concretas de internalización de los respectivos conocimientos, con un criterio superador de las regulaciones coyunturales, de modo que el alumno pueda, ante cualquier reforma legislativa, tener herramientas suficientes para adaptar los conocimientos adquiridos, al eventual nuevo diseño normativo, sin perder la sustancia de la institución.

\section{Conocimientos previos}
La asignatura no requiere contenidos técnicos jurídicos previos, más allá de los que pue
den ser parte integrante de la currícula de los planes de estudio de las escuela d
e nivel medio y del Ciclo Básico Comun.

\section{Contenidos Mínimos}
Los contenidos temáticos de esta asignatura, priorizan a la persona como centro de
protección jurídica, cuestión ésta que resulta altamente positiva en atención a los fines a que se halla orientada la labor universitaria en general y que fueran referidos precedentemente.
\\ \\
Así, luego de la necesaria referencia a temas estructurales del Derecho Privado y
Civil en particular, necesarios para que el alumno comprenda el campo del conoc
imientos que nos ha de ocupar, la consideración de la persona en tanto sujeto de der
echo, tanto desde lo individual, como integrando un ente colectivo, sus atributo
s, la protección jurídica de hechos relevantes como su nacimiento y su muerte o el rég
imen regulatorio de su obrar lícito, en tanto se proyecta jurídicamente, se puede de
cir que constituyen los contenidos mínimos a modo de cimientos del edificio que lu
ego se dedicará a terminar de construir, integrando los demás conocimientos propios
del resto de las asignaturas.
\\ \\
\section{Objetivos Generales}
El desarrollo de los contenidos de los distintos ejes temáticos, se endereza a sat
isfacer dos tipos de objetivos: a) genéricos y b) específicos. \\ \\
En este plano, teniendo en cuenta los fines que entiendo la Universidad está desti
nada a alcanzar, así como la renovación constante del conocimiento y los problemas a
que se deberán enfrentar los futuros egresados en el medio socio-económico al que s
e incorporarán, se persigue que los alumnos:\\
\begin{itemize}
\item Aprendan a aprender,
\item Aprendan a pensar,
\item Aprendan a investigar,
\item Aprendan a comprender un problema y a resolverlo,
\item Aprendan a interrelacionar conocimientos, tanto entre sí como con sus propias vivencias
\item Se reconozcan como responsables directos y protagonistas de su formación.
\end{itemize}

\section{Objetivos Específicos}
Como objetivos específicos, se califican a aquéllos que resultan corolario de los an
tedichos genéricos, cuando se los confronta con el contenido puntual de la asignat
ura. En este contexto, se procura que el alumno:
Vivencie la utilidad de sus conocimientos teóricos,
Valorice el análisis de los fallos jurisprudenciales como dato de la realidad,
Incorpore como elemento integrador de su aprendizaje, el derecho comparado,
Comprenda las instituciones jurídicas en su esencia, más allá de la forma con que aparezcan consagradas en un determinado momento de la historia,
Tome conciencia y se comprometa en el estudio de los temas cuya consideración se a
bre al futuro.


\chapter{Unidad 1}
\section{¿Que es el Derecho?:}
El hombre es un ser social y no puede vivir aislado, aparece así la vida en sociedad para que esta sea posible es menester que exista una reglamentación que indique lo que necesariamente no debe hacerse y aquellos que forzosamente debe realizarse.
\\
En otras palabras cierta parte de la actividad individual deja de depender del libre arbitrio de cada uno y queda sujetada a las normas que se dictan en comunidad, el problema de separar las actividades que deben reglamentarse de las que deben quedar libradas a las decisiones individuales es un problema difícil que a lo largo de la historia ha recibido diversas soluciones seg\'un el punto de vista podemos citar al iusnaturalismo y el positivismo.
\\
\textbf{Esas normas que regulan la convivencia de la vida en sociedad es lo que llamamos el derecho es decir el ordenamiento social.}   \\
Algunos autores a ese ordenamiento le agregan el adjetivo de justo, \textbf{Arauz castex dice que el derecho es la convivencia humana normativa mente pensada en función de justicia}, El jurista llambías dice que el derecho es el ordenamiento social justo y en general todos los autores coinciden en que ese ordenamiento de la vida en sociedad debe tender al valor de justicia.

\subsection{concepto}
\subsubsection{ distintas definiciones articuladas por la doctrina}

. Importancia del concepto. 
Derecho y Moral. Diferencias y afinidades.
Derecho subjetivo. Concepto, fundamento: teorías varias. Teorías negativas de los

\subsection{Derecho objetivo y subjetivo.}

la palabra derecho tiene una variedad de contenidos semánticos la podemos utilizar en expresiones tales como el derecho civil o "tengo derecho a algo"
\\
Podemos ver que en el primer caso se utiliza el término derecho refiriéndonos al derecho objetivo (norma agenda) es decir el derecho exterior a la persona en tanto en la segunda nos referimos a una prerrogativa es decir una facultad o prerrogativa.
\\
\textbf{llambías define el derecho subjetivo como la facultad reconocida por el ordenamiento jurídico para exigir de los demás un determinado comportamiento borda en tanto la define como la facultad de exigir a otro una determinada conducta}
\\
la noción de derecho subjetivo nos presenta tres características en primer lugar tenemos el derecho objetivo que constituye la norma por ejemplo el artículo 2506 del código civil dice, El dominio es el derecho real en virtud del cual una cosa se encuentra sometida a la voluntad y la acción de una persona, Esta persona por ejemplo puede ser la propietaria de un inmueble y como tal puede disponer del mismo.
\\
Otra característica es el deber que las demás personas tienen de respetar los derechos de un propietario notemos que se habla de deber y no de obligación mientras toda obligación es un deber no todo deber es una obligación técnicamente sólo se denominan obligaciones a los deberes de carácter relativo y de contenido patrimonial a los cuales corresponde un derecho subjetivo creditorio.
\\
La tercera consecuencia se produce cuando por ejemplo un tercero molesta al propietario de un bien y éste recurre a la justicia en busca de una sanción que será una indemnización por daños y perjuicios por lo tanto la consecuencia de la norma jurídica son facultad, deber y sanción









Derecho Natural y positivo: concepto. Ramas tradicionales y modernas
Derecho civil: evolución histórica y contenido actual.
Principios Generales del Derecho. Concepto. Su importancia actual.
La codificación. El método en el Código Civil y Comercial de la Naciòn.
\\ 


Derecho objetivo y derecho subjetivo
El derecho objetivo es lo que acabamos de definir es el ordenamiento social justo es decir el conjunto de normas que regulan la vida en sociedad

el derecho subjetivo es la facultad con la prerrogativa que me reconoce el ordenamiento jurídico para exigir de una persona un determinado comportamiento

También debemos recordar que la palabra derecho tiene una variedad de contenidos semánticos podemos utilizar la expresión en por ejemplo el derecho civil y también podemos expresar lo siguiente tengo derecho a en el primer caso estamos utilizando el término como derecho subjetivo es decir exterior a la persona y en el segundo caso estamos ante la presencia del derecho subjetivo interior a la persona
Nos podríamos preguntar si es válida la distinción creemos que sí la expresión derecho objetivo se refiere a la norma jurídica la que rige las relaciones interhumanas y el término derecho significa una regla una disposición o un sentido una ley en este sentido el derecho objetivo se refiere a por ejemplo el derecho de los propietarios el derecho del padre etcétera es la norma agendi del derecho romano

Y usamos el vocablo para señalar la facultad la actitud la posibilidad de cada uno para hacer o no hacer algo para exigir de otros un comportamiento a ésta le llamamos facultad agendi romano

Derechos subjetivos
ya envían lo define como la facultad reconocida por el ordenamiento jurídico para exigir de los demás un determinado comportamiento borda por otro lado no dice que los derechos subjetivos es la facultad de exigir de otro una determinada conducta

La naturaleza jurídica de los derechos subjetivos por ejemplo savini un jurista alemán  decía que los derechos subjetivos es el poder atribuido a una voluntad esto es el ordenamiento jurídico me da un poder una facultad me empodera para que uno pueda hacer algo sin embargo las críticas que reciben es que pone el acento el acento en la voluntad lo que puede llevar extremos como el abuso

Von hiering por otro lado decía que el derecho subjetivo se define por la utilidad que le brinda al hombre



\chapter{clase 2 - Fuentes del derecho.}
Son los distintos modos de expresi\'on	del derecho positivo, el objetivo es determinar cuales son las causas, sociales, politicas o morales, que dieron o\'igen a una norma y cuales son los modos de expresi\'on del derecho.
\section{¿Quién es el intérprete?}
Obviamente es el juez pero también todos los que trabajamos en el derecho lo interpretamos, es decir cuando a uno le presentan un caso y le cuentan los hechos, hacemos una interpretación dentro del marco jurídico, pero en este texto cuando nos referimos al intérprete nos referimos al intérprete último es decir el juez.
\\
Las fuentes las vamos a dividir en fuentes formales y fuentes materiales.
\section{Fuentes Formales}Son las que obligan al intérprete.   
\begin{itemize}
\item la ley
\item la costumbre
\item los fallos plenarios
\end{itemize}
\subsection{La ley}
Es la primera de las fuentes formales, evidentemente primero existió la costumbre antes de la ley en los pueblos primitivos, y a medida que fue avanzando el tiempo la costumbre quedó de lado y la ley cobró mayor importancia.
¿Pero qué es la ley? La ley, es la norma obligatoria, dictada por una autoridad competente que tiene carácter de generalidad y tiene coactividad, es una norma que tiene socialidad está dirigida a la generalidad de las personas y no está dirigida a casos particulares, si nosotros tomamos el código civil el artículo 4, dice que las leyes son obligatorias para todo ciudadano de la república ya sean nacionales o extranjeros, residentes domiciliados o transeúntes, sin perjuicio de lo dispuesto en leyes especiales.
\\
\\
Tiene carácter de coactividad, obliga a los destinatarios si bien el hombre es libre de no cumplirla, pero habrá una sanción y esa es la coactividad.\\
La sanción típica del derecho civil es la nulidad, cuando se celebra un acto jurídico que no cumple con los requisitos, su sanción será la nulidad de dicho acto, pero la característica de la ley es la coactividad es decir obliga si no, sería una mera exhortación legislativa sin peso.
\subsubsection{Clasificación de la ley}Se clasifica según distintos criterios, por ejemplo, podemos hablar de ley en sentido formal y en sentido material,
\subsubsection{ley en sentido formal y en sentido material}
La ley en sentido formal es la que emana del congreso es fácilmente identificable tiene un nombre y un número y además, siguió el procedimiento que dicta la Constitución nacional.
\\
\\
\textbf{Por ejemplo} ley ley 27551 ley de locaciones, tiene un nombre, un número y siguió el procedimiento que dicta la Constitución nacional, el poder ejecutivo la promulga, el congreso sanciona y después se publica para que el público la conozca.
\\
\\
Nosotros llamamos ley en sentido material a cualquier disposición que emana de una autoridad competente, siempre y cuando tenga carácter de generalidad.
\\
\\
\textbf{Por ejemplo} una resolución de la afip es ley pero no es ley del congreso, pero es ley en el sentido que es una norma jurídica  y obligatoria, general, emanada de una autoridad competente o sea es ley en sentido material siempre que tenga generalidad,el carácter de generalidad no es que se tiene que aplicar a todo el mundo sino a un número indeterminado de personas por ejemplo la ley de locaciones no afecta a personas que no estén involucrados en una relación locativa, sólo se aplica a los inquilinos y a los locatarios.
\\
\\
\textbf{Otro ejemplo}, si la ley dice que a partir de tal fecha se eleva el impuesto en el monotributo esta ley se dirige a todos los que están incluidos dentro del monotributo y no a los demás entonces es una ley en sentido material.
\\
\\
Muchas veces una ley lo es en los dos sentidos, por ejemplo una ley que emana del congreso con el procedimiento que la constitución dicta y tiene carácter de generalidad, surge de una autoridad competente, por ejemplo la ley de locaciones es ley en sentido formal, porque salió del congreso pero también es ley en sentido material porque se dirige a un número indeterminado de personas y surgió de una autoridad competente pero también puede pasar que no sea ley en los dos sentidos, por ejemplo si el congreso dijera que se va a levantar un monumento a una persona en algún lugar, va a ser una ley en sentido formal porque emana de una autoridad competente pero no es ley en sentido material.
\\
\\
La ley de presupuestos es una ley en sentido formal no en sentido material,
una ley que dijera concédase de por vida una pensión a una persona por ser viuda de tal también es una ley en el sentido formal pero no en sentido material.
\\
\textbf{No hay que confundir ley en sentido material con fuente material o ley en sentido formal con fuente formal la ley ya sea en sentido formal o material y siempre una fuente formal porque obliga al intérprete.}

\subsubsection{La imperativa y ley supletoria}

Nosotros dijimos que una característica de las leyes su coactividad su obligatoriedad la ley obliga pero a veces lo hace de manera diferente

La ley imperativa obliga con toda la fuerza de su imperatividad y por lo tanto los particulares no la pueden dejar de lado, no pueden decir esta norma no me conviene o no me gusta y no tenerla en cuenta para nuestra relación jurídica, es decir no pueden dejarla de lado.

Por ejemplo, la ley de locación dice que el plazo mínimo debe ser de 3 años, esa es una ley imperativa y las partes no podrían decir vamos a hacer nuestro contrato por un año.

\subsubsection{Ley supletoria}
Es supletoria de la voluntad de las partes si bien la norma establece una determinada cuestión los particulares podrían dejarla de lado y regirse de otra manera, nosotros podríamos preguntarnos dónde está la coactividad o la imperatividad de la ley en ese caso, la imperatividad sigue pero más liviana, pero continúa, porque si los particulares pueden dejarla de lado sólo es porque la ley se lo permite.

Si las partes no convienen nada distinto a la ley vigente en el contrato, no pasa nada simplemente se aplica lo que la norma dice es decir suple la voluntad inexpresada de las partes.

La ley supletoria se aplicarán en el sector en el que rige soberana la autonomía de la voluntad qué es la parte de obligaciones y contratos en ese sector es más fácil que existan normas supletorias.

Por ejemplo la norma que dice que los padres le deben alimentos a los hijos hasta la mayoría de edad es una norma imperativa los padres no podrían decir que no les gusta la norma y en tal caso solamente quieren hacerse cargo del menor hasta los 10 años.

Pero en el caso de las locaciones, el concepto de Mora si es una norma supletoria.


\subsubsection{Ley Prohibitiva y dispositiva} Otra clasificación en el sentido de la disposición. \\

Primitiva es aquella ley que impone la abstención de alguna conducta y la dispositiva es la que ordena hacer algo.

En tal Barrio por ejemplo no edificar a más de cierta altura es una norma prohibitiva y una norma dispositiva es la que obliga al deudor pagar al acreedor sus acreencias

¿Cómo se cumple una obligación de no hacer?
Por ejemplo, no edificar hasta cierta altura se cumple no haciendo exactamente lo que se pide es decir absteniéndose de tal conducta, y la ley dispositiva se cumple llevando a cabo tal conducta.

Otra clasificación de acuerdo a la sanción
la antigua clasificación romana decía que una norma podría ser perfecta imperfecta más que perfecta.

La sanción típica del derecho civil es la nulidad, un acto jurídico que no cumple los requisitos que la ley establece el mismo no va a producir los efectos que busca es decir se va a producir la nulidad.

En este caso una norma perfecta es aquella que prevé la sanción de nulidad por ejemplo el dolo que es uno de los vicios de la voluntad, el dolo es esa maquinación es astucia de engañar a otra persona para que se celebre un contrato que sin el dolo no lo habría hecho.

Por ejemplo en un contrato de compraventa el cliente se dispone a comprar un smart TV El vendedor lo engaña y le dice que tiene esta característica sin embargo no la tenía entonces la víctima cumple los requisitos para pedir la nulidad de dicho acto
\subsubsection{ley perfecta}
Aquella norma que prevé la nulidad de un acto es una ley perfecta porque prevé la sanción de nulidad.

\subsubsection{ley imperfecta} es aquella que no prevé sanción en caso de incumplimiento no sería una verdadera ley sería como una especie de pedido del legislador de que hagan tal cosa pero si no hay sanción no es ley por eso la llamaban imperfecta.

\subsubsection{ley menos que perfecta} es la que prevé una sanción distinta de la nulidad, como dijimos en el caso del ejemplo del televisor para pedir la nulidad debemos cumplir cuatro requisitos si no se cumplen no podemos pedir nulidad pero la víctima puede pedir daños y perjuicios, entonces la norma que dice que no puede pedir la nulidad pero puede pedir daños y perjuicios es una ley menos que perfecta no es perfecta porque no prevé la nulidad no es imperfecta porque tiene sanción.

\subsubsection{ley más que perfecta o pluscuamperfecta}
Contiene la sanción de nulidad pero además también contiene otras sanción es más que perfecta es decir la nulidad más otra en el caso del dolo tiene sanción pero lo mejor todos los perjuicios que sufrió la víctima no quedan subsanados con la nulidad por ejemplo si el cliente perdió la oferta de otra tienda por comprar en ese lugar entonces además de la nulidad puede pedir daños y perjuicios.

\subsubsection{¿Cuándo entra en vigencia una ley? Y cuando deja de tener vigencia}

Hay todo un procedimiento, hay un proyecto de ley se discute se debate en diputados senadores y finalmente se llega al acuerdo y las ley es sancionada por el congreso, luego va el poder ejecutivo y el poder ejecutivo la promulga tácita o expresamente luego debe ser dada a conocer. 

Si la ley lo dice, entra en vigencia cuando la ley lo diga pero generalmente la ley no lo dice entonces entra en vigencia a los 8 días después de su publicación en el boletín oficial sin contar el día de la publicación porque los plazos se cuentan de medianoche a medianoche.

Por ejemplo
Si yo recibo una notificacion hoy, y me dan 3 dias para hacer algo, el día de hoy no se cuenta porque el código dice que un día va de medianoche a medianoche.

\subsubsection{¿Y hasta cuándo tiene vigencia?}
Tiene vigencia hasta que se dicta otra ley que la derogue, porque no pueden existir dos leyes que traten la misma cuestión en sentidos diferentes, esa derogación puede ser expresa o tácita.
por ejemplo cuando entró en vigencia el nuevo código civil de la nación decía por ley 26994 Desia expresamente, deróguese las siguientes leyes y presentaba una lista de las mismas esa se trata de una derogación expresa.

La derogación tácita es cuando una ley trata una cuestión en sentido opuesto a otra entonces la ley posterior deroga ley anterior tácitamente.

\subsubsection{Jerarquía de leyes}
Es como una pirámide que en la punta está la Constitución nacional qué es la ley suprema pero también al mismo nivel están los tratados internacionales de los que nuestro país fue parte las luego las constituciones provinciales, luego leyes nacionales y las leyes provinciales, Luego decretos y ordenanzas municipales y resoluciones.

\subsection{La Costumbre}
Es otra de las fuentes formales, estamos hablando de la costumbre como uso social sino la costumbre como fuente productora de normas jurídicas, la costumbre es un comportamiento una conducta general llevada a cabo por todos los miembros de una comunidad es un comportamiento reiterado prolongado en el tiempo uniforme homogéneo, existe la convicción por parte del agente de qué es una obligación 
Es un uso social saludar y qué pasa si no lo hago no pasa nada, y qué pasa si no le dejo propina al mozo tampoco pasa nada.
la costumbre por otro lado es una productora de normas jurídicas y tiene un elemento subjetivo y otro objetivo

\subsubsection{Podemos hablar de tres tipos de costumbres}
\subsubsection{La costumbre secundum legem} o sea costumbre según la ley es la costumbre que está reconocida expresamente por el ordenamiento jurídico la ley la menciona la reconoce
por ejemplo en un contrato la ley dice que si no está fijada la fecha de cumplimiento de la obligación esta será la que indican los usos y costumbres del lugar.
\\
\\
El contrato de locación por ejemplo los usos y costumbres de la ciudad de Buenos Aires dicen que se pague por adelantado y del 1 al 5 de cada mes eso es una costumbre secundum legem.

\subsubsection{La costumbre praeter legem }es la costumbre más allá de la ley y es la costumbre que llena un vacío de la ley cuando la ley no regula un tema la costumbre viene a llenar ese vacío.
\\
\\
Por ejemplo el apellido de la mujer casada hoy en día es una costumbre que se dejó de lado pero antes existe una ley dictada en el año 1969 ley 18248 qué decía la ley del nombre que le quedó drogada cuando entró en vigencia el nuevo código civil y comercial, esta ley decía que la mujer debía agregar a su apellido el apellido del marido precedida por la preposición de ejemplo Pérez de González, recién en el niño 69 hubo una ley del apellido pero antes era un uso o costumbre que las mujeres hagan uso del apellido del marido pero no había ley que lo obligue, pero no era solamente una costumbre simple en una costumbre que debía cumplirse a tal punto que hubo casos de divorcio porque la mujer no llevaba el apellido del marido, en este caso el juez declaró que había injuria por parte de la mujer y uno se pregunta si no existía la ley de dónde sacó el juez las bases para falar de esta manera entonces podemos inferir que se trataba de una norma consuetudinaria como no preveía nada la costumbre viene a llenar los vacíos de la ley.

\subsubsection{La costumbre contra legem} es una costumbre que se forma contraria a los dictados de la ley.
De estos tres tipos de leyes nuestro código acepta los dos primeros
\subsection{Fallos plenarios}
Los fallos pueden ser comunes es decir la jurisprudencia como un o fallos plenarios esto es cuando se reúne la cámara en pleno los fallos comunes cuando uno apela el fallo en primera instancia iba a la cámara, se sortea la sala, los fallos plenarios son dictados por todas las salas en conjunto por los 39 camaristas.
\\
\\
Uno como abogado recurre a los fallos de segunda instancia, ¿qué pasa cuando un fallo es contradictorio con otro de la misma cámara? esto es cuando un fallo contradice a otro, en el mismo tema y la misma c\'amara, cuando existe un caso de sentencias contradictorias sobre un mismo tema, hay métodos de unificar la jurisprudencia y uno de los métodos es convocar a la cámara en pleno.\\
\\
\subsubsection{Cómo se hace}
La cámara podría auto provocar la reunión en pleno pero sino una de las partes puede interponer un recurso llamado recurso de inaplicabilidad de la ley qué es un recurso en el que no se discuten los hechos se discute el derecho o la aplicación de la norma es decir no se discute los hechos si se chocó el auto en tal esquina o si el acusado dijo tal cosa se discute el derecho es decir la aplicación de la norma al caso.
\\
\\
Mediante este recurso se puede solicitar la reunión de la cámara en pleno si el fallo contradice alguno que se dictó dentro de los 10 años anteriores por alguna sala de la cámara. Cuando la cámara se reúne en pleno está y decide cuál es la doctrina aplicable.
\\
\\
Un caso conocido fue el de el comprador en un remate que adquiriera un inmueble, la cuestión es quién debe pagar las expensas atrasadas el comprador oel rematado todo eso fue un tema de una reunión plenaria otro tema también fue el de los intereses aplicables, el fallo plenario obliga a todos los jueces de la cámara y a todos los jueces inferiores a aplicar la resolución o doctrina que surge del fallo.
\\ \\
Entonces un fallo plenario \textbf{¿es fuente formal o fuente material?} claramente es fuente formal porque obliga, tradicionalmente los fallos plenarios fueron vinculantes osea obligatorios pero una ley en la que se crearon las cámaras de casación (casar en francés quiere decir anular), y dijeron que iban a sacar la obligatoriedad de los fallos plenarios porque las reuniones plenarias iban a ser reemplazadas por los tribunales de casación esos tribunales nunca se crearon entonces la ley 27500 de enero del 2019 le restituyó a los fallos plenarios su obligatoriedad por lo tanto los fallos plenarios continúan siendo vinculantes es decir obligan a todos los camaristas a que cuando tengan un caso igual resuelvan de la misma forma esto es la jurisprudencia plenaria, entonces la jurisprudencia común es fuente material porque no obliga \textbf{pero persuada pero la jurisprudencia plenaria es fuente formal porque obliga}


\section{Fuentes Materiales}Son las que lo convencen o lo persuaden de cuál puede ser la solución justa, son todas las demás fuentes del derecho, entonces son 
\begin{itemize}
\item la jurisprudencia común
\item la doctrina
\item el derecho comparado
\item los principios generales del derecho
\item la equidad
\end{itemize}
Son todas fuentes materiales que no obliga al intérprete pero sí son tenidas en cuenta por este a la hora de decidir.

\subsection{La Jurisprudencia}
es la reunión de los fallos de distintos jueces que pueden orientar a un tercero cuando éstos son coincidentes, para que se forme jurisprudencia necesario que la resolución de los jueces se vayan reiterando un solo fallo no hace jurisprudencia, no sienta precedentes
un juez dicta sentencia de acuerdo a su real saber y entender por lo tanto los fallos de otros jueces sobre un mismo tema no lo obligan sin embargo si los fallos se van retirando el juez puede orientarse en base a ellos
por eso decimos que la jurisprudencia es una fuente material y no una fuente formal.
\\
\\
Las sentencias de los jueces tienen vistos, fallos y considerandos, en el visto dice por ejemplo, visto que Juan Pérez inició acciones contra Pedro González etcétera, por daños y perjuicios.
\\
Los considerandos son la parte medular del fallo, el juez da los motivos por los cuales falla en la forma que lo hace, la sentencia de un juez podemos decir que tiene una parte reflexiva y una parte autoritaria, la parte reflexiva está en los considerandos y la parte de autoritaria es cuando dice visto tal cosa resuelvo tal otra.

\subsubsection{Artículo 3 del código civil} dice que el juez debe resolver los casos que lleguen a su jurisdicción a través de una solución razonablemente fundada, esto es una solución adecuada y suficiente, adecuada al caso y suficiente en el sentido de que existen suficientes fundamentos para tal decisión, cuando una apela una sentencia lo que tiene que rebatir son los considerandos.


\subsubsection{Cómo está organizado el foro civil}
Primero están los juzgados de primera instancia, qué son 110, el juez dicta la sentencia y si alguna parte no está de acuerdo puede apelar, entonces el expediente se eleva a la cámara nacional de apelaciones en lo civil, la cámara tiene 13 Salas identificadas de la A a la M y cada una tiene tres jueces, ellos tienen que decidir si se mantiene o no la sentencia de primera instancia, los verbos que vamos a ver en un fallo de segunda instancia es revocar o confirmar.
\\
\\
Y la última instancia, es la Corte suprema, obviamente son muy pocos los casos que llegan, la Corte suprema tiene competencia original y completencia derivada, la competencia original son los casos en donde la corte tiene competencia original por ejemplo el conflicto entre dos provincias pero tiene una competencia derivada cuando llegan apelaciones puedes llegar a la corte cuando se pone en tela de juicio la inteligencia de interpretación de alguna garantía de la Constitución nacional


\subsection{La doctrina}
Es la opinión de los autores es una fuente claramente material, si recordamos lo que dijimos de los considerandos en los que el juez tiene que dar los fundamentos de su decisión, los textos son doctrina y estos ayudan a estudiar el derecho es decir a conocer la ciencia del derecho y le da al juez el material para fundar su fallo no es una fuente secundaria que no tenga importancia.
\\
\\
A veces cuando se publica un fallo también la editorial introducir algunos comentarios del fallo, así tenemos dos fuentes jurisprudencia y doctrina se convoca algún autor y se le pide que haga algún comentario de dicho fallo.
\\
\\

\subsection{Principios generales del derecho}
Son aquellos principios fundamentales que algunos autores dicen que provienen del derecho natural que está inscripto en la mente del hombre no hay una lista exhaustiva de principios generales del derecho pero sí podemos mencionar algunos.
\\
\\
Por ejemplo la buena fe o que no se puede dañar a otros la buena fe lealtad y la buena fe creencia el principio del enriquecimiento sin causa el principio de irretroactividad de la ley la indemnización de equidad el derecho comparado, es decir hay una serie de principios que están presentes en las mentes de los hombres y que se consideran también con fuente material del derecho.

\subsection{el derecho comparado} en particular es fuente material en alguna materia que no está regulada entonces el juez podría decir que en otras legislaciones este tema ha sido tratado de tal manera.
\subsection{la equidad }
En el caso de la equidad recordemos todas las definiciones de derecho apuntan al sentido de justicia pero a veces la justicia de la norma deja un gusto amargo por ejemplo cuando una persona no tiene discernimiento por un padecimiento mental la norma dice que no tiene que indemnizar a la víctima porque no se le puede reprochar nada porque no entiende ya que por medio de la ley de los actos involuntarios no se responden sin embargo no deja de ser cierto que nos quedamos como que no se hizo justicia pero vino la ley 17711 e introdujo la indemnización de equidad

la ley dice que es cierto que no se puede reprochar nada a un demente porque es autor de un acto involuntario pero tampoco podemos dejar pasar el daño producido entonces ahí viene la equidad como decía Aristóteles a solucionarlos efectos de la justicia rigurosamente legal laaplicación de la norma a lo mejor es legal pero no es equitativa entonces viene la ley 17711 a corregirlo es una fuente material

los hechos se deben probar pero el derecho no sé yo digo fundir hecho en la ley de educación es no tengo que demostrar nada

\textbf{La costumbre se prueba o no se prueba}
Quizá un poco la costumbre en un juicio tengo que probarlo o no bueno si decimos que es un hecho lo tenemos que probar cidio que es una norma jurídica o sea es derecho entonces lo tengo que probar, pero siendo la costumbre una fuente formal no la tengo que demostrar
\chapter{fuentes del libro ameal}
Fuentes del derecho
Clasificación de las fuentes
Fuentes formales
Fuentes materiales

La ley la Constitución nacional los tratados de derechos humanos los usos prácticas y costumbres

La ley los casos deben ser resueltos conforma un sistema de fuentes pero en este sistema la principal fuente es la ley ya que de lo contrario aparecerían sentencias que no aplican la ley o se apartan de ella sin declarar su inconstitucionalidad siendo esto una decisión contra legem
Tiene dos acepciones
Ley en sentido amplio o material
Ley en sentido estricto o formal

Caracteres de la ley
Es obligatoria y coactiva
Generalidad
Origen público

La Constitución nacional
Los tratados de derechos humanos a los que se han adherido
Los usos prácticas y costumbres

La costumbre

Elementos y carácteres
Elemento objetivo o material la costumbre 
Para qué se de este elemento, debe reunir los siguientes carácteres

Ser uniforme ser constante largo uso generalidad y publicidad

Elemento subjetivo psicológico o espiritual se da cuando existe la firme creencia por parte de la comunidad de que lo que está haciendo es una necesidad jurídica y por lo tanto obligatoria

Tipos de costumbres

\includegraphics[scale=1]{./img/ley1.png}



\chapter{clase 3 - Irretroactividad de la ley}
El principio indica que la ley se aplica hacia el futuro y no hacia el pasado, el fundamento de este principio es la seguridad jurídica porque si uno adquiere un derecho cumpliendo todos los requisitos que la ley exige para ser titular del mismo y luego a posteriori viniera una ley más exigente que modifica esos requisitos, entonces no habría derecho alguno que se pudieran adquirir, por eso la tranquilidad de las personas que adquieren derechos se funda en la existencia de este principio.
\\
\textbf{¿A quién va dirigido este principio?, ¿A legislador o a los jueces?}
\\
A legislador claramente no, porque puede dictar leyes retroactivas lo que no puede hacer es mediante la sanción de una ley, violar derechos que están amparados en la constitución nacional.
El principio va dirigido al juez que cuando está por decidir sobre un caso sabe que tiene que aplicar la norma hacia el futuro y no hacia el pasado.
\\
\textbf{¿Cuál es el fundamento de este principio?}
Vélez sarsfield se fundaba en la \textbf{teoría de los derechos adquiridos}, esta teoría dice que si yo cumplo con todos los pasos que la ley exige para adquirir un derecho, el mismo ya forma parte de mi patrimonio y no puede venir una ley posterior a arrebatarmelo, en cambio decía Vélez que si podía afectar derechos en expectativa es decir, una esperanza de derecho que necesita la acción fecundante del tiempo para convertirse en algo tangible.
\\
\textbf{Ejemplo}
supongamos un ciudadano de 65 años que cumple todos los requisitos para jubilarse el anses le otorga dicho derecho y el ciudadano comienza a gozar de la jubilación, a posteriori se dicta una ley que eleva la edad jubilatoria, no podemos decirle al ciudadano que vuelva a trabajar porque ya no cumple los requisitos, en cambio si es un ciudadano que aún no cumplió los 65 años y se dicta la ley que eleva la edad jubilatoria, esa esperanza de derechos s\'i se ve afectada por la nueva norma, y no podemos decir que se esté afectando un derecho adquirido.
\\ \\
Vélez tenía organizado esto en los artículos 4044 y 4045, en el primero se refería los derechos adquiridos y en el segundo a los derechos en expectativa, él decía que la ley nueva s\'i afectaba derechos en expectativa o facultades que no se hubiesen ejercido.
\\
\\
En el código de Vélez había dos casos en el que sí se podía aplicar la ley retroactivamente uno era el orden público esto es que nadie tenía derechos irrevocablemente adquiridos contra una ley de orden público por lo tanto la ley podría afectar un derecho adquirido.
\\ \\
Definamos qué entendemos por \textbf{orden público}, son todas aquellas normas que regulan instituciones importantes para la existencia de las sociedades por ejemplo la familia la propiedad etc.
\\ \\
Otra excepción era la \textbf{ley aclaratoria}, a veces hay leyes que son oscuras o ambiguas y es necesario que se dicte otra ley para aclararla entonces la ley aclaratoria y la ley oscura forman un conjunto de una única ley tácita, en este caso la ley aclaratoria siempre vuelve a modificar o aclarar la ley anterior.
\\ \\
La teoría de los derechos adquiridos no era el único fundamento para explicar la existencia de este principio también está \textbf{la teoría de los hechos cumplidos}, de Rubíer un autor francés, esta teoría dice que cuando un hecho ocurrió en el pasado y produjo todas sus consecuencias es decir se agotó y se terminó, no puede venir una ley nueva a querer abarcarlo dentro de sus previsiones, ya está terminado el hecho esto se parece al concepto de llambías de consumo jurídico es decir ya se consumieron los efectos y la ley nueva no puede aplicarse a ese supuesto.
\\ \\
La ley 17711 modificó el código de Vélez, derogó los artículos 4044 y 4045 y el artículo 3 lo redactó de otra manera pero, sigue ratificando la existencia de este principio diciendo \textbf{toda ley es irretroactiva sea o no de orden público}, y dice específicamente que el legislador puede dictar leyes retroactivas siempre que no afecten derechos amparados en garantías constitucionales.
\\ \\
\textbf{¿Qué es lo que dice este artículo 3 que luego se convertiría en el artículo 7 del nuevo código?}
Que a partir de su entrada en vigencia las leyes se aplican a las consecuencias de las relaciones y situaciones jurídicas existentes.
\\
\textbf{Una relación jurídica} es una vinculación entre personas esencialmente variable destinada a durar poco tiempo, por ejemplo la relación entre las personas a través de un contrato.
\\
\textbf{Una situación jurídica}, en cambio, es de mayor permanencia por ejemplo la situación familiar o la relación de padre a hijo.

\subsection{Ley imperativa y ley supletoria}

\textbf{La ley imperativa} es aquella a la cual los particulares no pueden dejar de lado por qué se refieren a temas importantes para la sociedad y los particulares no pueden regirse de otra manera que no sea la que la norma establece.
\\
En cambio \textbf{la ley supletoria} está librada a la voluntad de las partes es decir si no les parece adecuada para su contrato dicha norma puede dejarse de lado y regirse de otra manera.
\\
\textbf{Por ejemplo}, la situación de mora es el retraso en el cumplimiento de una obligación, que es importante porque a partir de ahí se empiezan a calcular intereses, entonces una de las partes solicita que se lo interpele es decir que se le avisé por medio legal, antes de entrar en mora eso lo pueden hacer las partes pero si no dicen nada, la ley supletoria "suple" la voluntad no expresada de las partes.
\\
\textbf{¿Es distinta la aplicación de la ley cuando es supletoria o imperativa?}
Cuando la ley es imperativa y entre en vigencia se aplica a las consecuencias de relaciones y situaciones jurídicas existentes, si hubo un contrato que se celebró antes de la entrada en vigencia de la ley imperativa pero aún continúa produciendo sus efectos, en este caso a este tramo no cumplido del contrato, se aplica la ley nueva, eso no es retroactividad, eso es efecto inmediato de la ley nueva.
\\  \\
Hay contratos que se caracterizan por ser de cumplimiento instantáneo por ejemplo la compraventa pero hay otros contratos que duran a lo largo del tiempo, son de ejecución diferida, por ejemplo la locación, que su plazo mínimo es de 2 años otro ejemplo es un préstamo o mutuo.
\\
\textbf{Por ejemplo} los particulares celebran un contrato de préstamo a pagar en 12 meses a un 10\% mensual a los 6 meses de la celebración de dicho contrato, se sanciona una ley nueva imperativa que dice que no se puede cobrar más del 5\% de interés, a las 6 cuotas que faltan pagar se aplicará un interés del 5\%, y no es así en el caso de las 6 cuotas ya pagadas, si le exigimos la devolución del 5\% al acreedor por la entrada en vigencia de esta ley, eso sí sería retroactividad en cambio a esto se llama efecto inmediato de la ley imperativa.
\\
\textbf{Ley supletoria}
Las partes quisieron regularse de una determinada manera entonces, o bien cambiaron la forma de regularse \textbf{y lo expresaron en el contrato} o bien aceptaron lo que decía el código, si viniera una ley nueva a cambiar el contenido de ese contrato, sería como mover las bases del mismo, y las partes por otro lado, si hubieran sabido que vendría una ley nueva a cambiar las reglas de juego, tal vez no hubieran celebrado dicho contrato o se hubieran regulado de otra manera.
\\ \\
\textbf{¿Qué ocurre cuando entra una ley nueva supletoria y hay un contrato en ejecución? }
en ese caso no se aplica la ley nueva, se aplica la ley vieja, esto es una ley que ha sido derogada sigue aplicándose a un contrato vigente, es una situación particular, como si hubiese un muerto que aún trabajé, en el caso jurídico esta ley muerta (en el sentido de que ha sido derogada) continúa trabajando para el cumplimiento de dicho contrato, este efecto se llama \textbf{efecto ultractivo} prolongado o diferido de la ley supletoria.
\\
Esta regulación que surgió de la ley 17711 y que modificó el artículo 3 del código de Vélez, 
hizo una modificación sustancial  cuando se refiere a una ley supletoria, continúa diciendo que no entra en vigencia, si existe un contrato aún vigente a menos que la ley nueva sea más beneficiosa para el consumidor, es decir que cuando la ley nueva sea más favorable al consumidor (parte m\'as d\'ebil) la ley vieja ya no tendrá efectos en el contrato vigente.





\chapter{Ejercicio abusivo del derecho}

Teoría del abuso del derecho en el código civil Comienza en el capítulo 3 
Buena fe en el artículo 9 abuso del derecho artículo 10 abuso de posición dominante artículo 11 abuso del derecho individual en perjuicio del ambiente en el artículo 14 

Estos principios son utilizados en todas las ramas del derecho privado, por eso se introdujo en el código civil en el capítulo 3 con el título el ejercicio de los derechos de los ciudadanos 

El principio de buena fe
buena fe objetiva
Se aplica el ejercicio de todos los derechos y a cualquier relación jurídica, este principio dice que los derechos deben ser ejercidos de buena fe, es decir, con lealtad o buena fe objetiva, 
buena fe subjetiva
en tanto la creencia de que se está actuando correctamente es la buena fe subjetiva.


La ley concede derechos subjetivos
En la clase pasada vimos la diferencia entre el derecho objetivo y el derecho subjetivo Que recordamos que el derecho objetivo era el ordenamiento jurídico y el derecho subjetivo son las prerrogativas de los ciudadanos.

Por ejemplo, el propietario de un bien tiene derecho a disponer del mismo, el acreedor tiene derecho a cobrar su deuda y esto está garantizado jurídicamente. Cuando la ley otorga estos derechos lo hace teniendo en cuenta un fin determinado, dicho fin tiene que ser el buscado por el legislador cuando pensó la norma, es decir el acreedor en el ejercicio de su derecho subjetivo no puede despegarse de los fines que se tuvieron en mente al momento de sancionar la ley que otorga dicho derecho.


Es decir los derechos subjetivos se deben ejercer dentro de determinados límites como ser los fines que tuvo el ordenamiento jurídico al reconocer el derecho subjetivo en cuestión, los límites impuestos por la buena fe la moral y las buenas costumbres, es decir no afectar al ambiente y a los derechos de incidencia colectiva en general.

Si una persona se aparta de estos fines y límites que impone la justicia estamos en presencia de un ejercicio abusivo de un derecho subjetivo, por eso se dice que el ejercicio de los derechos subjetivos es relativo y no absoluto, porque son limitados  Es decir lo tenemos que ejercer dentro de ciertos límites. 

Otros síntomas de estar en presencia de un abuso en el ejercicio del derecho pueden ser el ejercer el derecho sin necesidad sin interés o de una manera razonable o excesiva o también ejercer el derecho con dolor con lo con la sola intención de perjudicar.


Surgimiento e Historia
El Código de Vélez sarsfield no tenía en cuenta los derechos subjetivos y no tenía en cuenta el ejercicio el derecho es decir un derecho no podía ejercerse abusivamente porque sería un una contradicción.

El ejercicio regular de un derecho propio o el cumplimiento de una obligación legal no pueden constituir como ilícito ningún acto, la ley no ampara el ejercicio abusivo de los 
se considera tal, los que contrarían los fines del ordenamiento jurídico o el que excede los límites impuestos por la buena fe la moral y las buenas costumbres, el abuso del derecho es el ejercicio de un derecho en forma irrestricta sin límites y sin respeto por el otro, consiste en un límite al ejercicio del derecho subjetivo.

El viejo código civil de Vélez existía el artículo 1071 que decía que el ejercicio regular de un derecho propio o el cumplimiento de una obligación legal no puede constituir como ilícito ningún acto eso significa que no se podía considerar un derecho como abusivo,y mucho menos se enseñaba la teoría del abuso del derecho.

La reforma constitucional de 1949 en el gobierno del General Peron, si contemplo la teoría del abuso del derecho, lo hizo en el artículo 35 Y decía algo así como los derechos y garantías reconocidos por esta constitución no podrán ser alterados por las leyes que reglamenten su ejercicio Pero tampoco amparan a ningún habitante de la nación en perjuicio, detrimento o menoscabo de otro, el ejercicio abusivo de los derechos que perjudican a la comunidad  o que lleven a cualquier forma de explotación del hombre por el hombre configuran delitos que serán castigados por las leyes.

Con este artículo del año 49 se introdujo la teoría del derecho en la legislación Argentina Hubo fallos o sentencias que hicieron uso de esa teoría, es decir la jurisprudencia empezó a hacer uso de esa teoría y con la reforma de la ley 17711 se aceptaba esta teoría, dicha ley modifica el artículo 1071 escribiéndose de la siguiente forma, El ejercicio regular de un derecho propio una obligación jurídica no puede constituir como ilícito ningún acto y continúa diciendo la ley no ampara el ejercicio abusivo de ningún derecho, se considera tal el que contraríen los fines que aquélla tuvo en mira al reconocerlo o los que contraríen la una fue la moral y las buenas costumbres. 


Habíamos dicho que el código de Vélez sarsfield no tenía en consideración la teoría del abuso del derecho y los fundamentos para esto eran los siguientes Si el estado se constituye en juez del abuso muy pronto se convertiría en juez del uso y toda idea de derecho subjetivo se perdería, sin embargo el artículo 2503 de la parte de derechos reales dice Algunos ejemplos de cómo no ejercer los derechos.

Recordemos que la constitución del 49 fue abolida en el golpe militar del 55 Al derogarse esta Constitución la situación volvió a la del Antiguo régimen Es decir no había ninguna norma en la que apoyarse para fundamentar el ejercicio abusivo de un derecho Aparece en la constitución del 49 el artículo 35 Como fundamento para fallar en tal sentido.
Luego del derrocamiento del General perón se deja sin efecto la constitución del 49  Y volvemos Estado en el que no había ninguna Norma que hablar al respecto Luego tuvimos que esperar hasta el año 1968 con la ley 17711 para que se incorporarse definitivamente esta teoría al código civil 




Entonces el titular del derecho podría ejercerlo pero tenía que hacer dentro de los límites razonables que imponían la moral la buena fe y las buenas costumbres Es decir de una manera regular.


Criterio subjetivo

En el artículo 10 del actual código civil tenemos,

Para observar el abuso de un derecho podemos hablar de distintos criterios  Criterio subjetivo  Círculo Con la idea de culpa Funda el acto abusivo En la intención de perjudicar a otro 	Sin que su actuación lo beneficie Este criterio está vinculado con el dolo definiendo dolo como la intención de dañar Y que se ejerce el dicho derecho si no tener ningún beneficio Entonces el dolo la culpa y la falta de interés serían los tres elementos de este criterio .
El dolo 
es la intención de perjudicar Es decir cuando ejerce el derecho de tal manera que no teniendo beneficio del mismo perjudica a un tercero Secretario fue criticado en su momento cuando apareció por qué En realidad no había muchos casos en los que se gestión derecho con la intención de perjudicar Por eso se incorporó la culpa Es decir no solamente ejercicio abusivo del derecho cuando el que lo ejerce tiene intención de perjudicar Sino también cuando ejerce el derecho con culpa Entendiendo la culpa como el actuar negligente de alguna persona Es decir el famoso hacer de menos Es decir quien ejerce el derecho Él no busca el daño pero se produce el daño De igual manera Por la falta de diligencia Es decir es un hacer de menos no poner los cinco sentidos en lo que está haciendo 

la culpa
Pero también la culpa es hacer demás La imprudencia también es tu aspecto de la culpa es aquella persona que hace La negligencia es hacer de menos y la imprudencia es hacer demás

falta de interés
Y el otro criterio subjetivo también Es la falta de interés Podríamos decir porque uno ejerce el derecho subjetivo De esta manera si no tengo ningún interés en el sentido de que no tengo ningún beneficio Y además estoy dañando un tercero Entonces el juez Se preguntará si no tiene interés de hacerlo en esta manera y aún así lo está haciendo Y además está dañando un tercero por qué lo hace Porque precisamente lo que busca es perjudicar.


Concepción objetiva del abuso del derecho Habría abuso del derecho cuando se ejerce En oposición a los fines Económicos y sociales Que esperaron la ley que los contiene.

Existe también una posición en mixta En donde se combinan los dos criterios anteriores 


El criterio objetivo o finalista es muy utilizado porque

Todos los derechos se conceden para algo cuando el legislador reconoce un derecho en una norma lo hace con una norma con una finalidad determinada

Por ejemplo para que se le otorga a los padres lo que antes se llamaba patria potestad y hoy se llama responsabilidad parental para que se le drogaba los padres ese cúmulo de derechos y deberes de los padres respecto de los hijos menores de edad se les da esas facultades para someterlos para que ejercen un poder sobre ellos??
Ciertamente no se da esa responsabilidad de los padres para que puedan educarlos bajo sus criterios el legislador entonces al momento de conceder esos derechos tuvo en mira el bienestar del menor ahora sí esos derechos se utilizan de otra manera o para hacer otra cosa como por ejemplo o como un modo de sojuzgar al hijo.

ejemplo para que le da el legislador el derecho al acreedor para cobrar sus acreencias y poder embargar al deudor algún bien de su propiedad en caso de incumplimiento porque el legislador le da ese derecho
Para vengarse no evidentemente xeloda para que pueda satisfacer sus acreencias entonces todos los derechos tienen una finalidad que la ley tiene en mira al momento de concederse cuando el ejercicio del derecho se desvía de ese fin estamos en presencia de un ejercicio abusivo del derecho

El nuevo código sancionado no le da tanta importancia al legislador sino al ordenamiento jurídico ese es el criterio finalista

Otro criterio es el de la moral y las buenas costumbres el código nos da una doble directiva por un lado nos da el criterio finalista y por otro el de la moral y las buenas costumbres


¿Qué decir cuando hay ejercicio abusivo?
Es muy difícil aceptar que los derechos no se deben ejercer abusivamente el problema radica en saber cuándo es abusivo el ejercicio tenemos que ver la pauta que nos dé el código

Será un juicio quiero utilizar la teoría del abuso del derecho a lo tengo que encuadrar en lo que dice el código por ejemplo diciendo que el ejercicio de tal derecho va en contra del artículo 10 del código civil o va contra la moral y las buenas costumbres

otros casos de uso del derecho el código civil y comercial de la nación contempla el abuso en el ejercicio de un derecho por parte de su titular pero además se incluyen otros casos de abuso por ejemplo el artículo 10 habla de la situación jurídica
esto es un significativo desequilibrio entre los derechos y las obligaciones de las partes en perjuicio del consumidor artículo 1119 y 1120 del código civil

El artículo 1119 dice está dentro del capítulo 4 que dice cláusulas abusivas y dice sin perjuicio de lo dispuesto en las leyes oficiales es abusiva la cláusula que habiendo sido o no negociado individualmente tiene por objeto o por efecto provocar un desequilibrio significativo entre los derechos y obligaciones de las partes en perjuicio del consumidor

El artículo 1120 dice se considera que existe una situación jurídica abusiva cuando el mismo resultado se alcanza a través de la predisposición de una pluralidad de actos jurídicos conexos

Ejemplo voy a sacar un paquete de viaje de turismo la empresa tiene derecho a recibir el pago del viaje pero a su vez me dice que yo debo pagar una suma de gastos adicionales es una cláusula dentro del contrato si seguimos leyendo el contrato hay otra cláusula que dice lo siguiente habiéndose iniciado el viaje la empresa no tiene la obligación de explicarme en que utiliza dicha suma

¿Esa cláusula es abusiva o no?
Es abusiva porque porque si el dinero es utilizado por ejemplo para excursiones entonces está claro en qué se utiliza pero si no se utiliza deberían devolverlo
entonces en dicho contrato las dos cláusulas conjuntamente hacen una situación abusiva es decir hay un desequilibrio entre las partes en perjuicio del consumidor

Eso también lo podemos ver en la medicina prepaga en los contratos bancarios estar van a ver en la materia contratos y obligaciones

artículo 11 del código civil de la nación abuso de posición dominante lo dispuesto en los artículos anteriores se aplica cuando se abusa de una posición dominante en el mercado sin perjuicio de las disposiciones específicas contempladas en las leyes especiales ejemplo consiste en que una empresa que tiene una posición de dominio en el mercado abusa de esa situación de superioridad sobre sus competidores o clientes para imponer por ejemplo excesivas condiciones comerciales perdiendo Independencia comercial

esto será regularmente en los casos de monopolio u oligopolio por ejemplo cuando existe un único vendedor puede imponer el precio que quiera

abuso del derecho individual que afecta el ambiente o derechos de incidencia colectiva en general artículo 14 del código civil el ejercicio de un derecho individual no puede afectar el funcionamiento ni la sustentabilidad de los ecosistemas de la flora la fauna el agua el paisaje etcétera artículo 240 del código civil

el código civil establece numerosas normas que hacen referencia a casos de abuso del derecho por ejemplo en materia de contratos en artículo 988 y 989 cláusulas abusivas en materias de adhesión en materia de contratos de consumo artículos
1117 y 1118
Las cláusulas que desnaturalizan las obligaciones del periódico al ente

Hay que prestar atención para no confundir un ejercicio abusivo del derecho con un acto ilícito

Cuando hablamos del ejercicio abusivo de un derecho estamos en presencia de un derecho reconocido por la ley cuyo ejercicio es con alguno de los tres criterios mencionados el dolo la culpa etcétera es decir lo está ejerciendo fuera de los límites que el ordenamiento jurídico tuvo en mira al momento de concedernos pero existe el derecho y aquí hacemos mención a la contradicción que algunos autores plantean porque si existe el derecho no puede ser abuso y si existe un abuso no puede ser de un derecho pero hay que recordar que es el ejercicio lo abusivo y no el derecho

otra cosa es que se trata de un acto ilícito aquí no hay ningún derecho y ejerce una actividad contraria al ordenamiento jurídico a título de dolo o de culpa eso es un acto ilícito y la característica fundamental es que provoca daño

En el caso recordado en el de la quema de pastizales existe ejercicio abusivo o es un acto ilícito evidentemente no hay ningún derecho en este caso pero si puede ser el caso de que existe el derecho de ejercer alguna actividad que se ejerció imprudentemente





Efectos

el juez debe ordenar lo necesario para evitar los efectos del ejercicio abusivo es decir procurar la reposición al Estado de hecho anterior y fijar una indemnización

La buena fe creencia
este número de los casos en que la buena fe creencia el código lo tiene en cuenta

Por ejemplo el caso del heredero aparente no en cualquier caso se aplica por ejemplo yo creo que soy heredero de tal persona y me comporto como tal pero yo no sabía que no estaba incluida entre los herederos ese es un caso de la buena fe creencia

La buena fe en el tema de los vicios se reitera cuando una persona desconocía la existencia de algún vicio el orden jurídico lo reconoce

por ejemplo vicios de actos jurídicos como el fraude y la simulación

Para el caso particular de la teoría del abuso del derecho la buena fe que se debe tener en cuenta es la buena fe objetiva o la buena fe lealtad es decir yo tengo que ejercer mi derecho de buena fe con respeto a los derechos de la otra parte sino se desvirtúa el derecho si cada uno interpreta el ejercicio a su modo

Ejemplo un señor que vive en una propiedad horizontal y tiene una mascota el reglamento de copropiedad de el edificio establecía que no se podían tener mascotas evidentemente este es una norma que el señor conocía porque al momento de firmar la compra o la locación de un departamento uno se somete a los dictados de dicho reglamento en este caso no se permiten a mascotas sin embargo el copropietario tenía una entonces el consorcio lo intima a que se deshaga de la mascota el señor no hace caso y el consorcio le inicia un juicio por incumplimiento del reglamento exigiendo al juez que obligará a dicho señor a que cumple con el reglamento

Señor invoca la teoría del abuso del derecho y cuando uno invoca esta teoría está reconociendo que hay un derecho en la otra parte pero lo está ejerciendo de forma abusiva respecto a su persona porque si tienen varios copropietarios que también estaban en esa situación y a ellos no se estaba exigiendo el cumplimiento del reglamento

Otro ejemplo un vecino estaba haciendo una construcción y construye sobre unos pocos centímetros sobre el terreno de su vecino bueno el vecino se dio cuenta pero no dijo nada y espero a que termine en la construcción y le exigió a que destruyera lo construido aquí existe un ejercicio abusivo del derecho porque el damnificado vio desde un inicio que se estaba construyendo en su terreno y no dijo nada porque espero a que se termine toda la construcció??

Otro caso es el tema de trabar embargos el acreedor tiene derecho a cobrar y también tiene derecho a utilizar las herramientas que el ordenamiento jurídico realidad para llegar ese fin
Entonces se intima al deudor a que pague el deudor no paga entonces se inicia una acción judicial y trabajamos un embargo entendiendo éste como una medida cautelar que inmoviliza un bien en el patrimonio del deudor si es un bien mueble no lo puede mover y ciencia un bien inmueble no puede disponer del mismo entonces traigo un embargo para asegurar su derecho de cobrar sus acreencias pero suponiendo que la deuda es ínfima y el acreedor trabó un embargo a todos los bienes del deudor pero tengo uno no podría embargar algún bien proporcional a la deuda entonces esto se considera un ejercicio abusivo del derecho de embargo, alomejor el abogado no averiguo bien y ejerció abusivamente el derecho de embargo negligentemente

cómo hacían los jueces antes de que se incorporará la doctrina del de abuso del derecho en el código??

En general recurrían a otras normas como no había normas en muchos casos recurrieron a la legislación comparada que recordemos que no es obligatoria pero sí es una fuente material también recurrieron al artículo 953 del código de Vélez y ese artículo decía cuando se refería al objeto del acto jurídico y decía claramente como no debe ser el objeto del acto jurídico era una característica distintiva en el que definía en lugar de definirlo como debe ser lo definía como no debe ser

Y decía que el objeto del acto jurídico no debía ser contrario a la fe la buena fe la moral y las buenas costumbres entonces cuando un juez recibía una invocación a la teoría del abuso un derecho antes de la introducción de esta teoría en el código entonces el juez decía qué tal ejercicio y va contra la moral la buena fe y las buenas costumbres, tampoco se había incorporado a la lesión que es un vicio que vamos a ver más adelante









\section{Teoría clásica}
\subsection{Definición de variables}
\subsection{Pruebas y refutaciones}
\section{Hipótesis}
\chapter{Resultados}
\section{Simulación de resultados}
\subsection{Suposiciones}
\subsection{Modelos}
\section{Resultados preliminares}
\section{Resultados postprocesados}
\subsection{Valores atípicos}
\subsection{Correlaciones}
\chapter{Conclusiones}
\end{document}
